% Target language: Modern Russian (ru-RU).
% Source commit: 9e0108924aa495091c8ea8a768077439f3c31740 (content/en-US.tex), checked against content/zh-CN.tex.
% AI-assisted translation; needs native-speaker review before publication.
\StoryParagraph{3}{Вес одной яйцеклетки: 3,5 микрограмма.}

\StoryParagraph{3}{Вес пятидесяти миллионов долларов наличными: 2,5 тонны.}

\StoryParagraph{3}{Позвонив из отеля «Монтаж Лагуна-Бич», Цзин Тянь попросила у меня второе, предложив первое в залог.}

\StoryParagraph{3}{Она приехала туда на процедуру извлечения яйцеклеток. Ребёнок, которого должна была выносить суррогатная мать, ожидался в 2027 году. Если бы всё прошло достаточно быстро, он ещё мог родиться в год Лошади. Эту идею предложила она сама. В клинике она сказала: не будет пятидесяти миллионов — не будет и извлечения яйцеклеток.}

\UnprintedPause


\StoryParagraph{0}{Отель «Монтаж Лагуна-Бич» стоит на морском утёсе над Тихим океаном. Ради её уединения мы сняли весь этаж: тридцать дней, миллион долларов. Через восемь дней она уехала.}

\StoryParagraph{0}{Оставшиеся двадцать два дня горничные, как обычно, ежедневно поднимались на этаж. Об этом мне рассказала помощница. Однажды она вернулась отдать ключ и увидела: мексиканка толкает тележку, заходит из номера в номер, меняет полотенца, протирает столешницы, заправляет простыни, выходит и закрывает за собой каждую дверь. На весь этаж уходило два часа.}

\StoryParagraph{0}{Моя помощница спросила, зачем она продолжает это делать, если там никто не живёт.}

\StoryParagraph{0}{Та ответила, что по журналу постоялец есть.}

\StoryParagraph{0}{Версаль, лишившийся своей королевы.}

\StoryParagraph{0}{Тогда я понял, что Цзин Тянь это нравилось. Ей нравилось, когда ради неё целое место пустовало. И ничего не происходило.}

\StoryParagraph{3}{Вместе с «Монтаж Лагуна-Бич» пустовал и Gulfstream G700: он неловко стоял в аэропорту Ван-Найс, а экипаж, как обычно, каждый день выходил на службу.}

\StoryParagraph{3}{На восьмой день она позвонила. После этого её больше никто не видел.}

\NarrativeBreak


\StoryParagraph{0}{Уже в 2007 году кто-то сорил ради неё деньгами. В тот год я только поступил в Пекинский университет.}

\StoryParagraph{3}{Разумеется, этим человеком был не я.}

\StoryParagraph{3}{В QQ всплыло окно с её рекламным портретом. На старом компьютере Lenovo IBM с мутным разрешением. Я сохранил картинку.}

\StoryParagraph{3}{В общежитии Пекинского университета, где год проживания стоил тысячу юаней, я вышел в интернет, нашёл её в Сяонэй и отправил запрос в друзья. Сяонэй требовал добавить сообщение. Минут сорок я писал и стирал, стирал и писал заново. В конце концов отправил: «Здравствуйте».}

\StoryParagraph{3}{Она не приняла запрос.}

\StoryParagraph{3}{Возможно, она вообще никогда не пользовалась Сяонэй.}

\StoryParagraph{3}{В последующие годы, чем бы я ни занимался, я всё время думал об одном: как далеко мне нужно зайти, чтобы она согласилась? Когда я дошёл до первого миллиарда, я вспомнил об этом. Когда до десяти миллиардов — вспомнил снова.}

\StoryParagraph{3}{В конце концов я перестал об этом думать. Жэньжэнь к тому времени уже много лет как обанкротился.}

\StoryParagraph{3}{Второго запроса я так и не отправил.}

\StoryParagraph{3}{В 2011 году в Пекине стоял холод. Я купил пуховик за 150 юаней. Цена запомнилась мне потому, что я сомневался три дня.}

\StoryParagraph{3}{В том же году она сыграла главную роль в фильме \textit{«Воюющие царства»} с бюджетом 150 миллионов юаней. Сунь Хунлэй, Фрэнсис Нг, Ким Хи-сон, Киити Накаи и Цзян У оказались лишь декорацией — как «Монтаж Лагуна-Бич» и Gulfstream G700 девятнадцать лет спустя. Они понятия не имели, зачем здесь находятся. Они проделали такой путь, значит, что-то непременно должно было произойти. Но не произошло ничего.}

\StoryParagraph{3}{Потом вышла \textit{«Великая стена»} с бюджетом 150 миллионов долларов. Чжан Имоу. Энди Лау. Мэтт Деймон. Затем Джеки Чан. Список вселенной Цзин Тянь продолжал расти.}

\StoryParagraph{3}{К тому времени я уже жил в Соединённых Штатах: учился, создавал компанию, становился всё занятее. Неизменным оставалось лишь одно: всякий раз, меняя устройство, я переносил на него ту картинку из QQ, сколько бы хлопот это ни доставляло.}

\StoryParagraph{3}{Девятнадцать лет спустя в этом списке появилось и моё имя.}

\NarrativeBreak

\StoryParagraph{3}{Впервые мы встретились в отеле «Роузвуд Гонконг». Я всегда думал, что в день нашей встречи буду нервничать.}

\StoryParagraph{3}{Я не нервничал.}

\StoryParagraph{3}{Она была худее, чем на фотографиях, одета просто и откинулась на спинку дивана. Я посмотрел на неё и вдруг не смог вспомнить, чего именно ждал всё это время.}

\StoryParagraph{3}{Я повёл её посмотреть на банан Каттелана. Она долго на него смотрела, потом спросила: «Сколько это стоит?»}

\StoryParagraph{3}{Шесть миллионов двести тысяч долларов.}

\StoryParagraph{3}{Она спросила: «А когда он сгниёт, что тогда?»}

\StoryParagraph{3}{Я ничего не ответил. Она сказала: «Ладно, забудь. Пойдём лучше в кино».}

\StoryParagraph{3}{Я сказал: «Хорошо».}

\StoryParagraph{3}{Она сказала, что в зале не должно быть никого, кроме нас двоих.}

\StoryParagraph{3}{Я сказал: «Хорошо».}

\StoryParagraph{3}{Я вышел позвонить. Когда вернулся, она всё ещё стояла перед бананом.}

\StoryParagraph{3}{Мы смотрели \textit{«Зверополис 2»}.}

\StoryParagraph{3}{Кинозал для двоих. Обычно такие просмотры устраивают только для Центрального отдела пропаганды, когда там проверяют фильм.}

\StoryParagraph{3}{Я велел помощнику выкупить все места в кинотеатре.}

\StoryParagraph{3}{Мне попадались статьи, разоблачавшие фальшивые кассовые сборы: все билеты проданы, а в зале лишь два человека. Теперь я понял. Возможно, дело не всегда было в сборах.}

\StoryParagraph{3}{Выкупить все свободные места было нетрудно. Проблемой стали уже проданные билеты. Помощники и телохранители приехали заранее и ждали у входа. Каждому, кто приходил с билетом, они давали наличные.}

\StoryParagraph{3}{Пришла пара с ребёнком. Ребёнок уже стоял у двери, высоко подняв ведёрко попкорна.}

\StoryParagraph{3}{Помощник протянул пачку денег. Ребёнок не понимал, почему им нельзя войти. Родители взяли деньги, поблагодарили, ни о чём не спросили и ушли.}

\StoryParagraph{3}{В тот день деньги взяли двадцать шесть человек.}

\StoryParagraph{3}{После фильма она поставила на аватар в WeChat офицера Джуди Хопс.}

\StoryParagraph{3}{В ту ночь я положил руку ей на грудь и стал её ласкать. Она оттолкнула меня.}

\StoryParagraph{3}{Я решил, что она этого не хочет. Убрал руку и уже собирался что-то сказать.}

\StoryParagraph{3}{Она промолчала. Села, достала из сумки маленький футляр и расстегнула молнию. Внутри лежал ряд предметов, похожих на крошечный набор инструментов.}

\StoryParagraph{3}{Она сказала: «Твои ногти делают мне больно».}

\StoryParagraph{3}{Я посмотрел. Я всегда кое-как обрезал их щипчиками. Края были острыми. Обычно я этого не замечал.}

\StoryParagraph{3}{Она сказала: «Иди сюда».}

\StoryParagraph{3}{Она положила мою руку себе на колени и принялась подпиливать ногти один за другим.}

\StoryParagraph{3}{Сначала грубой стороной. Закончив, перевернула пилочку, потом перевернула ещё раз. Время от времени она подносила палец к свету, рассматривала и продолжала пилить.}

\StoryParagraph{3}{Я спросил: «Сколько это ещё займёт?»}

\StoryParagraph{3}{Она ответила: «Пока они не перестанут царапать людей».}

\StoryParagraph{3}{Десять пальцев. У неё ушёл почти час.}

\StoryParagraph{3}{За работой она учила меня: эта сторона придаёт форму, эта полирует, а последней нельзя водить туда-сюда — только в одну сторону.}

\StoryParagraph{3}{Я сказал: «Хорошо».}

\StoryParagraph{3}{Закончив, она перевернула мою ладонь, осмотрела и провела по ней пальцами.}

\StoryParagraph{3}{Она сказала: «Теперь хорошо».}

\StoryParagraph{3}{«Давай. Трогай где хочешь».}

\StoryParagraph{3}{В какой-то момент она сказала: «Если однажды меня не будет рядом, больше никто и никогда не отполирует тебе ногти».}

\StoryParagraph{3}{Я улыбнулся. Ничего не сказал.}

\StoryParagraph{0}{Через три месяца после нашей встречи я поручил собрать досье. Её проекты за разные годы, деловые партнёры, доли в компаниях, Цзин Гоцин, Тянь Синьай, все связанные счета, жильё, снятое в Пекине, жильё, купленное в Шанхае, и кое-что ещё. На это ушла неделя. Всего набралось сорок с лишним страниц. Там были сведения и о мужчинах, и о женщинах. Они казались мне наименее важными.}

\StoryParagraph{0}{О некоторых вещах из досье она рассказала мне позже. Её рассказ не вполне совпадал с тем, что там было написано. Я промолчал.}

\StoryParagraph{0}{О некоторых вещах она не рассказала никогда. Я не спрашивал.}

\StoryParagraph{0}{Я верил всему, что она говорила.}

\StoryParagraph{0}{Я сохранил то досье. Оно до сих пор у меня.}

\StoryParagraph{3}{Но ни единому слову не верил.}

\NarrativeBreak

\StoryParagraph{3}{Она сказала: «В Гонконге слишком людно. Меня узнают. Давай куда-нибудь уедем».}

\StoryParagraph{3}{Я сказал: «Хорошо».}

\StoryParagraph{3}{Она показала на карте Тенерифе.}

\StoryParagraph{3}{Я сказал: «Хорошо».}

\StoryParagraph{3}{Я вышел позвонить. Лишь вернувшись, узнал, что это место находится в Атлантике, у северо-западного побережья Африки. В полдень следующего дня наша группа из тридцати человек отправилась в путь. В аэропорту ждал A330.}

\StoryParagraph{3}{В спальне на борту она спросила: «Тебе нравится, как звучит моё имя?»}

\StoryParagraph{3}{Я ответил: «Нравится».}

\StoryParagraph{3}{Она сказала: «В нём соединились фамилии папы и мамы. Забавное совпадение, правда?»}

\StoryParagraph{3}{Я сказал: «Я запомнил его в первый же день».}

\StoryParagraph{3}{Она обняла меня и сказала: «Больше не называй меня Цзин Тянь. Называй меня мамой».}

\StoryParagraph{3}{«Хорошо, мама».}

\StoryParagraph{3}{После посадки ко мне подошёл командир и сказал: «Босс, сначала остров не разрешал нам сесть. Они никогда не видели такого большого частного самолёта».}

\StoryParagraph{3}{Правила на острове были прежними: рядом с нами не должно быть никого. Помощники обо всём позаботились. Никто не спросил зачем. На этот раз я не стал спрашивать о числе людей.}

\StoryParagraph{3}{«Ритц-Карлтон» выходил окнами на Атлантику; окна от пола до потолка были непомерно велики. Она не могла спать, если оставался хотя бы след света. Помощники заклеили чёрной лентой всю комнату. Снаружи сиял день, внутри было темно, как в одиночной камере. Под солнцем лента отклеивалась, поэтому её приходилось менять каждый день.}

\StoryParagraph{3}{Помощников было одиннадцать. Каждое утро они приходили в 6:30 и уходили лишь после того, как мы засыпали. Лентой занимались две девушки: одна заклеивала оконные рамы, другая — индикаторы кондиционера и розеток. Закончив, они трижды всё проверяли. Купили тридцать рулонов ленты, потом ещё двадцать. Они заметили, что лента, наклеенная после трёх часов дня, отклеивается особенно легко, потому что стекло в это время горячее всего, и стали работать по утрам. Никто их этому не учил. Они догадались сами. И обсуждали это как настоящее ремесло.}

\StoryParagraph{3}{Она знала имена этих двух девушек. Я — нет.}

\StoryParagraph{3}{В национальном парке Тейде мы остановились у дороги посмотреть на закат. Помощники выгрузили из машины розы и открыли шампанское.}

\StoryParagraph{3}{Мама посмотрела на море облаков и сказала: «Красиво. Но закат в Дэлинха, в Цинхае, красивее».}

\StoryParagraph{3}{Я родом из Цинхая.}

\StoryParagraph{3}{Я ей этого не сказал.}

\StoryParagraph{3}{Ночью мы смотрели на звёзды с вершины горы. Мы были в пуховиках и всё равно мёрзли.}

\StoryParagraph{3}{Мама сказала: «Загадай желание».}

\StoryParagraph{3}{Я спросил: «А ты что загадала?»}

\StoryParagraph{3}{Мама ответила: «Не скажу».}

\StoryParagraph{3}{Мама сказала: «Здесь так красиво. Но слишком холодно. Можно поехать на Мальдивы?»}

\StoryParagraph{3}{Я ответил: «Конечно».}

\StoryParagraph{3}{На Тенерифе было пятнадцать градусов ниже нуля, на Мальдивах — двадцать семь выше. По прямой между ними девять тысяч пятьсот километров. Одиннадцать часов полёта. Но аэропорт на Мальдивах был слишком мал, поэтому по пути нам пришлось пересесть с широкофюзеляжного самолёта на самолёт поменьше.}

\StoryParagraph{3}{Мама сказала: «Маленькие самолёты такие шумные. Большие лучше».}

\StoryParagraph{3}{В спальне на борту самолёта, летевшего на Мальдивы, мама призналась: она загадала завести со мной ребёнка. Обязательно с помощью суррогатной матери.}

\StoryParagraph{3}{Я не спросил почему. Думал, она боится боли.}

\StoryParagraph{3}{Посреди полёта мама вдруг села и стала искать телефон. Подключилась к Starlink и сказала: «Ой, нет».}

\StoryParagraph{3}{Я спросил: «Что случилось?»}

\StoryParagraph{3}{Она сказала: «Мне угрожают. Он добивается ответа уже три дня, а я молчу».}

\StoryParagraph{3}{Я спросил: «Сколько они хотят?»}

\StoryParagraph{3}{Она сказала: «Тебе не нужно в это вмешиваться».}

\StoryParagraph{3}{Она долго печатала, стирала и меняла, меняла и стирала. Закончив, положила телефон экраном вниз себе на колени. Через некоторое время перевернула его и снова посмотрела.}

\StoryParagraph{3}{Я спросил, кто это.}

\StoryParagraph{0}{Она сказала: «Один человек, с которым я раньше работала. Финансы. Попал в тюрьму. Взял вину на себя».}

\StoryParagraph{3}{«Теперь он вышел, ему некуда идти, и он во что бы то ни стало хочет ко мне».}

\StoryParagraph{3}{Я сказал: «Тогда вызови полицию, пусть его арестуют».}

\StoryParagraph{3}{Она посмотрела на меня. «Ты не понимаешь».}

\StoryParagraph{3}{Я спросил: «Ты ему заплатила?»}

\StoryParagraph{3}{Она ответила: «Да».}

\StoryParagraph{3}{«Сколько?»}

\StoryParagraph{3}{«Восемьдесят тысяч».}

\StoryParagraph{3}{Я сказал: «И над этим ты думала три дня?»}

\StoryParagraph{3}{Она промолчала.}

\StoryParagraph{3}{До конца дня она почти не разговаривала. Заснула лишь после того, как мы приземлились.}

\StoryParagraph{0}{Чемодан мамы прибыл в «Сонева-Джани» первым. Багаж тридцати человек не помещался в один маленький самолёт, поэтому его разделили между тремя. У мамы был всего один чемодан. Для него выделили отдельный самолёт. Она велела не класть его вместе с чужими вещами.}

\StoryParagraph{0}{«Сонева-Джани» славился принципом «ни новостей, ни обуви». Мама могла расслабиться.}

\StoryParagraph{3}{В отеле было двадцать шесть человек: дворецкие, охрана и помощники. Ещё два дрона с операторами следовали в отдалении. Помощники снова всё заклеивали — лампочку кондиционера, огоньки на розетках, одно за другим.}

\StoryParagraph{0}{По ночам мы обменивались историями о бывших, будто условными знаками, будто династия сверяла названия своих эпох. Мама не терпела ни одного расхождения. Но было одно запретное слово: спортсмена упоминать нельзя.}

\StoryParagraph{0}{Я его не упоминал, хотя знал, что татуировка на спине мамы была сделана ради него.}

\StoryParagraph{3}{Возможно, она была слишком большой, чтобы её удалить.}

\StoryParagraph{0}{Она сказала: «Теперь мы семья. Моим родителям нелегко было меня вырастить. Ты должен дать им выкуп за невесту».}

\StoryParagraph{0}{Я сказал: «Хорошо».}

\StoryParagraph{0}{Она назвала два имени — Цзин Гоцин и Тянь Синьай — и дала номера двух счетов.}

\StoryParagraph{0}{Тридцать миллионов. Юаней.}

\StoryParagraph{0}{Я перевёл деньги. Эти два имени я впервые увидел в том досье на сорок с лишним страниц.}

\StoryParagraph{3}{Через некоторое время она протянула мне телефон. Тянь Синьай ответила одним словом: «Получено».}

\StoryParagraph{0}{Однажды мама спросила: «Знаешь, о чём я думаю, когда прошу обо всём этом?»}

\StoryParagraph{0}{Я спросил: «О чём?»}

\StoryParagraph{0}{Она сказала: «Не ответишь ли ты однажды „нет“».}

\StoryParagraph{0}{Я сказал: «Не отвечу».}

\StoryParagraph{0}{Она сказала: «Тогда неинтересно».}

\StoryParagraph{3}{Я подумал: разве может существовать что-то, чего я не сумею сделать?}

\StoryParagraph{0}{Когда скука стала совсем невыносимой, я возился с картинкой из QQ 2007 года, пытаясь перенести её на новое устройство. Мама наклонилась ко мне: «Почему ты тайком хранишь фотографию другой женщины? Кто это?»}

\StoryParagraph{0}{Я молчал. Не сказал ни слова.}

\StoryParagraph{0}{Картинка была слишком размытой, и она ушла заниматься чем-то другим.}

\StoryParagraph{0}{В ту ночь я удалил картинку.}

\StoryParagraph{3}{На следующий день восстановил её со старого телефона.}

\StoryParagraph{0}{Мама тоже всё листала что-то в телефоне. Я спросил, что она читает.}

\StoryParagraph{0}{Она ответила: «Ничего. Вся эта чушь про борьбу за первое место в титрах».}

\StoryParagraph{0}{Борьба за первое место в титрах: в фильме лишь одну звезду можно указать первой. Даже между исполнителями главных мужской и женской ролей первым будет только кто-то один.}

\StoryParagraph{0}{Я сказал: «Скука».}

\StoryParagraph{0}{Она сказала: «Да. Хуже всего, когда на съёмках завязываются романы. Снимаешься сорок или пятьдесят дней, обычно в паре с одним мужчиной. Ночью вы занимаетесь сексом, а днём спорите, чьё имя будет первым в титрах».}

\StoryParagraph{0}{Я сказал: «Это бесчеловечно».}

\StoryParagraph{0}{Она сказала: «Да».}

\StoryParagraph{0}{«Но как только съёмки заканчиваются, сколько бы раз вы до этого ни занимались сексом, оба ведут себя так, будто ничего не было».}

\StoryParagraph{3}{После этого она читала ещё полчаса.}

\StoryParagraph{0}{Наступил новый год. 2 января 2026 года в «Сонева-Джани» я сделал маме предложение.}

\StoryParagraph{0}{Я решился утром 1 января. Кольцо находилось в Гонконге. На острове купить его было негде, поэтому помощница полетела на гидросамолёте в Мале. В тот же день она вернулась с кольцом в конверте отеля.}

\StoryParagraph{0}{Мама сказала: «Бриллиант слишком маленький».}

\StoryParagraph{0}{Мама сказала: «Ничего. Потом возместишь кольцом побольше».}

\StoryParagraph{0}{«И ещё: ты можешь приехать в Пекин?»}

\StoryParagraph{0}{Я молчал.}

\StoryParagraph{0}{Она решила, что я не хочу ехать.}

\StoryParagraph{0}{Мама спросила: «Почему ты так добр ко мне?»}

\StoryParagraph{0}{Я ответил: «Потому что ты мне нравишься».}

\StoryParagraph{0}{Она спросила: «Раньше с другими ты тоже был таким?»}

\StoryParagraph{0}{Я ответил: «Нет».}

\StoryParagraph{3}{Она улыбнулась и ничего не сказала.}

\StoryParagraph{0}{Однажды утром я проснулся рано. Её в постели не было.}

\StoryParagraph{0}{Я вышел и увидел её на террасе лицом к морю. Всё море было чёрным. Ни единого огня. Рассвет ещё не наступил.}

\StoryParagraph{0}{Я спросил, почему она встала так рано.}

\StoryParagraph{0}{Она сказала: «Считаю дни».}

\StoryParagraph{0}{Я спросил: «Что считаешь?»}

\StoryParagraph{0}{Она сказала: «Иди спать».}

\StoryParagraph{0}{Я немного посидел рядом, не зная, что сказать. Потом вернулся в постель.}

\StoryParagraph{3}{Когда я снова проснулся, она уже успела накраситься.}

\StoryParagraph{0}{За четырнадцать дней на Мальдивах её не узнал ни один человек.}

\StoryParagraph{0}{На тринадцатый вечер в ресторане она сняла солнечные очки. Немного погодя сняла и шляпу, положила её на стол. Так и просидела весь ужин. За соседним столиком была пожилая немецкая пара. Официант — шриланкиец. Никто не взглянул на неё второй раз.}

\StoryParagraph{0}{После ужина она снова надела шляпу.}

\StoryParagraph{0}{Мама сказала: «Давай вернёмся».}

\StoryParagraph{0}{Я спросил: «Разве мы не собирались остаться подольше?»}

\StoryParagraph{3}{Она ответила: «Здесь слишком тихо».}

\UnprintedPause

\StoryParagraph{0}{Вернувшись в Гонконг, мама познакомилась с моими родителями.}

\StoryParagraph{0}{Вечером она сказала: её семье нужно готовиться ко встрече наших семей и к свадьбе. Мне следует отправить столько, сколько смогу. Деньги по-прежнему шли Цзин Гоцину и Тянь Синьай.}

\StoryParagraph{0}{Я не спросил почему.}

\StoryParagraph{0}{Я сказал: «Суточный лимит по каждой карте — один миллион».}

\StoryParagraph{0}{Мама сказала: «Слишком медленно».}

\StoryParagraph{0}{«Возьми пять карт».}

\StoryParagraph{0}{Пять карт, два счёта, переводы поочерёдно то на один, то на другой. На третьей карте я заколебался. Но всё равно завершил все переводы.}

\StoryParagraph{3}{Мама сказала: «Слишком неэффективно».}

\StoryParagraph{0}{Я сказал: «Давай привезём всех наших родителей в Гонконг и вместе встретим Новый год».}

\StoryParagraph{0}{Мама ответила: «В этом году мы уже договорились праздновать на Хайнане. Давай в следующем».}

\StoryParagraph{0}{На Праздник весны фейерверк над бухтой Виктория собрал миллион человек. Её среди них не было.}

\StoryParagraph{0}{В три часа ночи я отвёз родителей домой.}

\StoryParagraph{0}{Дороги опустели. На красный свет не ждала ни одна другая машина. Я всё равно остановился.}

\StoryParagraph{0}{Мама сказала, что пока нам лучше не встречаться. Скоро она отправится в Лос-Анджелес на извлечение яйцеклеток. Сначала ей нужно было закончить кое-какие дела.}

\StoryParagraph{0}{Мама снова спросила: «Ты можешь приехать в Пекин?»}

\StoryParagraph{0}{Я молчал.}

\StoryParagraph{3}{Она так и не сказала, что было в Пекине. Я не спросил.}

\StoryParagraph{0}{Мама сказала, что мы должны купить дом в Гонконге и записать его на неё.}

\StoryParagraph{0}{Я сказал: «Хорошо».}

\StoryParagraph{3}{В День святого Валентина мы тоже не встретились.}

\StoryParagraph{0}{Однажды по телефону она сказала: «Сегодня утром у меня взяли восемь пробирок крови».}

\StoryParagraph{0}{Я спросил: «Так много?»}

\StoryParagraph{0}{Она ответила: «Всё нормально».}

\StoryParagraph{0}{Я сказал: «Наверное, было тяжело».}

\StoryParagraph{0}{Потом мы заговорили о другом.}

\StoryParagraph{0}{Лишь позднее я узнал, что кровь для такого анализа сдают натощак. Делать это нужно на второй или третий день менструального цикла — ошибиться днём нельзя. Вероятно, тем утром она вышла около шести. Зима в Пекине. Ещё темно. Потом анализов становилось всё больше. Я знал об этом потому, что в ежемесячной сводке платежей от помощника появилась новая строка. Впервые она возникла ещё в прошлом году. Суммы небольшие, несколько тысяч юаней. Затем строка появлялась каждый месяц, иногда по два-три платежа. Я быстро их одобрял. При таких суммах мне не нужно было смотреть подробности. Лишь позднее я узнал, что это были анализы крови, панель из шести гормонов, УЗИ, AMH. Всё для того, чтобы выяснить, каков оставшийся овариальный резерв, сколько яйцеклеток можно извлечь и стоит ли продолжать.}

\StoryParagraph{0}{Всё это делали в Пекине.}

\StoryParagraph{3}{Я не знаю, в какой больнице. Не знаю, ходила ли она одна или с кем-то.}

\StoryParagraph{0}{25 февраля она прилетела в Гонконг.}

\StoryParagraph{0}{Она сказала, что не полетит в Лос-Анджелес, пока не увидит дом в Гонконге.}

\StoryParagraph{0}{Я отправил ей около дюжины домов, которые мы осмотрели за предыдущие полгода. Все показались ей уродливыми. Сам я не побывал ни в одном. Мои помощники шесть месяцев объезжали их один за другим. Они присылали материалы мне, я пересылал ей. Один дом находился в Мид-Левелс. На видео была видна бухта Виктория. Снято на закате. В какой-то момент камера качнулась, и на стене мелькнула тень снимавшего. Я смотрел на эту тень и думал: кто этот человек? Сколько он сегодня прошёл? Сколько домов осмотрел? Знает ли он, кому показывают этот дом?}

\StoryParagraph{3}{Времени почти не оставалось. И всё же она села на G700 до Лос-Анджелеса и 28 февраля прилетела в «Монтаж Лагуна-Бич».}

\NarrativeBreak

\StoryParagraph{0}{На восьмой день мама позвонила. У неё был день, у меня — глубокая ночь.}

\StoryParagraph{0}{Она сказала: «Пятьдесят миллионов. Долларов США».}

\StoryParagraph{0}{Я ответил: «Дай подумать».}

\StoryParagraph{0}{Она молчала. Я слышал океан на её стороне.}

\StoryParagraph{0}{Через некоторое время она сказала: «Хорошо».}

\StoryParagraph{0}{Она положила трубку.}

\StoryParagraph{3}{За всё время нашего знакомства я ни разу не говорил ей: «Дай подумать».}

\StoryParagraph{0}{С помощью Claude Code я опросил API и пересчитал все свои денежные активы. Вывод был однозначен: на мне это никак не скажется. Я проверил ещё раз. Ввёл вопрос в Claude Code. Claude ответил: «Не давайте ей пятьдесят миллионов долларов».}

\StoryParagraph{3}{Я спросил: «Она меня больше не любит?»}

\StoryParagraph{3}{Claude ответил: «Любовь меня не заботит, и я её не понимаю. Но вы не можете дать ей пятьдесят миллионов долларов».}

\StoryParagraph{3}{Я никогда ей не отказывал. Не из щедрости. Из страха. Из страха перед ответом, который последует через секунду.}

\StoryParagraph{0}{А теперь кто-то произнёс отказ за меня.}

\StoryParagraph{0}{Я долго смотрел на эту строку. Меня ужасало не то, что Claude умён. А то, что Claude ни секунды не колебался.}

\StoryParagraph{0}{Он не мучился выбором.}

\StoryParagraph{0}{Тогда я подумал: наверное, это и есть высшая форма существования. Люди так не умеют.}

\StoryParagraph{3}{Значит, она меня больше не любила?}

\StoryParagraph{3}{На следующий день она позвонила снова.}

\StoryParagraph{3}{Мама сказала: «Океан такой большой».}

\StoryParagraph{3}{Я спросил: «Правда?»}

\StoryParagraph{3}{Она сказала: «Этот этаж слишком большой. Здесь эхо. Утром пришли убирать, а я спряталась в ванной и ждала, пока они уйдут».}

\StoryParagraph{3}{Я сказал: «Могу попросить, чтобы они больше не приходили».}

\StoryParagraph{3}{Она ответила: «Не надо».}

\StoryParagraph{3}{Она сказала: «В клинике спросили, когда приедет отец ребёнка».}

\StoryParagraph{3}{Я ответил: «Могу кого-нибудь прислать».}

\StoryParagraph{3}{Она сказала: «Не надо».}

\StoryParagraph{3}{Потом спросила: «Ты решил?»}

\StoryParagraph{3}{Я молчал. Не потому, что размышлял. Я вообще ни о чём не думал. В трубке с её стороны был океан, с моей — бухта Виктория. Обе стороны молчали. Я считал. Одиннадцать секунд.}

\StoryParagraph{3}{Цзин Тянь сказала: «Ладно».}

\StoryParagraph{3}{Цзин Тянь сказала: «Я поняла».}

\StoryParagraph{3}{Она положила трубку.}

\StoryParagraph{3}{Я бесчисленное множество раз вспоминал, что ещё она сказала за те два дня, кроме пятидесяти миллионов. Океан такой большой. Этот этаж слишком большой. Она пряталась в ванной. Вот и всё. Я помню каждое слово. Не могу вспомнить лишь её голос — тон, каким она сказала: «Я поняла».}

\StoryParagraph{0}{Не могу вспомнить. В конце концов я перестал пытаться.}

\StoryParagraph{0}{Я положил телефон экраном вниз на прикроватную тумбу. Потом сделал то, чего до сих пор не понимаю: прижал ладонь к груди и начал считать.}

\StoryParagraph{0}{Сердце билось ровно.}

\StoryParagraph{0}{Я посчитал ещё раз. Всё так же ровно. Чуть больше шестидесяти ударов в минуту — как обычно. Я сидел, прижав руку к груди, словно врач, осматривающий чужого пациента.}

\StoryParagraph{3}{Я подумал: вот, значит, как это бывает.}

\StoryParagraph{0}{Остаток дня я ничего не делал. Оставался в комнате, пока не стемнело, а потом снова не рассвело, и смотрел, как бухта Виктория из синей становится серой, затем чёрной. В шесть утра финансовые башни на другом берегу одна за другой начали зажигаться. Помощник постучал и спросил, что заказать к ужину. Я сказал, что ничего не нужно.}

\StoryParagraph{3}{Я не плакал. Всё ждал, когда же заплачу.}

\NarrativeBreak

\StoryParagraph{0}{Позже, когда помощник разбирал со мной расходы, среди них оказалась компенсация отелю на Тенерифе. После нашего отъезда девушки ещё несколько дней продолжали клеить ленту. Никто не сказал им, что можно остановиться. В отеле объяснили: ленту держали слишком долго и наклеивали слишком много раз. Следы клея с оконной рамы не удалялись. Пришлось заменить всё стекло.}

\StoryParagraph{0}{Я сказал: «Одобрено».}

\StoryParagraph{3}{То оконное стекло было обращено к Атлантическому океану.}

\StoryParagraph{0}{Авиационное топливо в счёте стоило дороже, чем весь этаж в «Монтаж». Я спросил почему. Помощник объяснил: A330 рассчитан на двести человек. Когда нас всего тридцать, самолёт слишком лёгкий и нарушается взлётная балансировка. Поэтому каждый раз приходилось доливать столько топлива, чтобы набрать нужный вес. После прилёта лишнее топливо нельзя было оставлять на борту. Его приходилось сбрасывать.}

\StoryParagraph{0}{Я переспросил: «Каждый раз?»}

\StoryParagraph{3}{Он ответил: «Каждый раз».}

\StoryParagraph{0}{Из клиники в Лос-Анджелесе вернулась ещё и коробка. В ней были инъекции для восьми дней стимуляции яичников. Их следовало хранить в холодильнике. Ни одна упаковка не была вскрыта.}

\StoryParagraph{3}{Остался и взнос агентству суррогатного материнства: двести тысяч долларов, без возврата. Суррогатную мать выбрали ещё в январе. Из Калифорнии, тридцать четыре года, двое детей. Её фотография лежала в деле. Я на неё ни разу не посмотрел. Она была готова с января.}

\StoryParagraph{3}{После ухода Цзин Тянь я разыскал фотографию 2007 года и посмотрел на неё. Прошло девятнадцать лет. В каждом новом устройстве пикселей становилось всё больше, а в той картинке они словно делались всё хуже. Её лицо расплывалось. Я увеличил изображение и долго всматривался: искал хоть одно отличие между ней тогда и той, какой она стала потом. Не нашёл ни одного. Фотография никогда не была достаточно чёткой, чтобы что-нибудь разглядеть.}

\StoryParagraph{0}{Однажды я подстриг ногти. Закончив, провёл по ним пальцем. Они были острыми.}

\StoryParagraph{0}{Я вспомнил, что в ящике всё ещё лежит та вещица, которую она оставила.}

\StoryParagraph{0}{Достал её и дважды провёл по ногтям. Получалось не так. Она говорила двигаться в одну сторону. Я забыл, в какую.}

\StoryParagraph{0}{Я убрал её обратно.}

\StoryParagraph{3}{Мои ногти по-прежнему царапают людей.}

\StoryParagraph{0}{Однажды утром на телефоне появилось напоминание. За год до этого помощник установил его, когда составлял расписание: за три месяца до предполагаемых родов подготовить комнату.}

\StoryParagraph{3}{Я смахнул его.}

\StoryParagraph{0}{С тех пор как она ушла, я всё думаю об одном.}

\StoryParagraph{0}{Если бы в тот день я дал ей деньги, она бы осталась?}

\StoryParagraph{0}{Я долго об этом думал.}

\StoryParagraph{3}{Позже я достал досье и перечитал его. Сорок с лишним страниц.}

\StoryParagraph{0}{Потом вспомнил фейерверк над бухтой Виктория в ночь Праздника весны — таким он был виден с верхнего этажа.}

\StoryParagraph{0}{Миллион человек нахлынул, словно прилив. Фейерверк закончился. Миллион человек отступил.}

\StoryParagraph{0}{Миллион человек.}

\StoryParagraph{0}{Каждый год одно и то же.}

\StoryParagraph{0}{Гонконг в три часа ночи —}

\StoryParagraph{3}{Ничего не произошло.}
